\section{Universality}
\label{sec:universality}

We assemble the pieces of Section~\ref{sec:automaton} into a register machine and state the main theorem.

\subsection{The register on the automaton}

A register is a chain of memory switches, laid as track and switch cells in the cell graph. Its value is the
number of units currently set. The locomotive increments it by entering the chain, running over the set units
to the first free one, setting it, and returning. It decrements by clearing the last set unit and returning.
When the register is zero the locomotive finds no unit to clear and returns by the zero track. Each of these is
a path through tracks and the three switches, every step of which is governed by the rules of
Section~\ref{sec:automaton}. By Lemmas~\ref{lem:direction} to~\ref{lem:switches} the locomotive follows the
intended path and leaves the switch settings in the intended state, so the register holds and updates its
value correctly.

\subsection{The program}

A register machine is a finite list of instructions, each either an increment of a register followed by a jump
to another instruction, or a decrement that branches on whether the register was zero. We lay each instruction
as a sub-circuit that drives the locomotive to the register it touches, performs the operation, reads the zero
track where the instruction is a decrement, and routes the locomotive to the next instruction. A loop or a jump
is track that returns. The whole program is a finite pattern, and the registers extend it outward in a fixed
periodic way, so the configuration is ultimately periodic, as weak universality allows.

We have implemented the register machine at the level of the railway and checked that it computes. A
two-register program for multiplication, run by the single locomotive, returns the product of its inputs, in
agreement with a reference register-machine interpreter \cite{pollard2026vibe}. This confirms that the railway
pieces, once correct, compose into a working universal machine, independently of any tiling.

\subsection{The theorem}

\begin{theorem}[Uniform weak universality]
\label{thm:main}
There is one cellular automaton, defined by rules local in the cell-adjacency graph and independent of any
tiling, which is weakly universal on every regular hyperbolic tiling. On a given tiling, a register machine is
run by taking as the initial configuration the track network of the machine embedded in the cell graph of the
tiling, with track and switch roles and labels as in Section~\ref{sec:setting}.
\end{theorem}

The automaton is the fixed rule of Section~\ref{sec:automaton}. By Lemmas~\ref{lem:direction}
and~\ref{lem:motion} it realizes a track on which a head-and-trailing-cell locomotive advances in one
direction, in any tiling. By Lemma~\ref{lem:switches} it realizes the fixed, flip-flop, and memory switches. A
crossing is realized as in Section~\ref{sec:setting}, by the four-way piece in the plane or by routing through
the extra dimension in space. These are exactly the pieces Proposition~\ref{prop:railway} requires, so the
automaton simulates an arbitrary register machine, which is universal by Minsky \cite{minsky1967}. The
embedding of a particular machine is the initial configuration, which is ultimately periodic, so the
universality is weak in the usual sense. None of the rules referred to the tiling, so the same automaton serves
for every regular tiling. \qed

\subsection{Scope of the claim}

Two remarks bound the theorem honestly. First, the construction needs the track network to be embeddable in
the cell graph of the tiling. Every finite graph embeds in a hyperbolic tiling because of the exponential room,
and crossings are resolved by the crossing piece, so the embedding always exists for a tiling with cells of
degree at least three. The degenerate cases are the few tilings too small to route a switch, which do not arise
among the hyperbolic regular tilings.

Second, the automaton is weakly universal, not strongly universal. The initial configuration is infinite, the
unbounded registers, though ultimately periodic. This matches the setting of the tiling-specific automata it
generalizes, which are also weakly universal, with the single exception of the strongly universal dodecagrid
automaton \cite{margenstern2021dodeca5}, which starts from a finite seed by a separate device. A
finite-seed, track-building variant of the present construction is possible and is taken up in companion work.
